\scanpage{10}
% archive foliation: 9
dire que $L_*\simeq L^0_*\otimes_{\mathbf Z}k$, où $L^0_*$
est un complexe de $\mathbf Z$-modules \unsure{$\frac12$ simpliciaux}
spécial.

\uline{Proposition 2} Les $k$-modules \inserted{$\frac12$ simpliciaux} spéciaux
(resp. str. spéciaux) sont stables par
produits tensoriels [et par les op.
tensorielles $\Lambda^i$, $\Gamma^i$, $\operatorname{Sym}^i$?]
\margin{\unsure{faux}}
\margin{[\ ] faux}

Le cas strict. spécial se ramène
au cas spécial, en se plaçant sur $\mathbf Z$\unsure{\ldots}

1° Produit tensoriel. On se réduit
à prouver que les modules
\[
\left\{\begin{array}{ll}
\Lambda^i\Phi_*^{\unsure{*}}\otimes\Lambda^j\Phi_* &(i,j\geq1)\\
\Lambda^i\Psi_*\otimes\Lambda^j\Psi_* &(i,j\geq0)\\
\Lambda^i\Phi_*\otimes\Lambda^j\Psi_* &(i\geq1,\ j\geq0)
\end{array}\right.
\]
sont spéciaux. Or ils sont à composantes
libres de \unsure{t.f.|r.f.}, et par Künneth-Eilenberg-
Zilber leurs $\pi_i$ sont itou, il reste à
montrer qu'il n'y a qu'un nb fini
de facteurs $M_i$, $N_j$ $\neq0$. \deleted{\unsure{On va}} \inserted{Ce résultat}
\deleted{\unsure{prouver}} résulte du

\deleted{\uline{Lemme} Sous les conditions du cor 2 de
prop 1 sont encore équivalentes celles
suivantes (en supposant \deleted{\unsure{Spec $k$ connexe}}, \inserted{$k$ connexe}
de sorte que $\mu_i=\operatorname{rg}M_i$ et $\nu_j=\operatorname{rg}N_j$ est
défini)

a) $\operatorname{rg}_kL_n$ est un polynôme en $n$
pour $n$ grand

b) $\exists$}
